\documentclass[11pt,a4paper]{article}
\usepackage[utf8]{inputenc}
\usepackage[T1]{fontenc}
\usepackage{amsfonts}
\usepackage[french]{babel}
\frenchbsetup{StandardLists=true} % à inclure si on utilise \usepackage[francais]{babel}
\usepackage{enumitem}
\usepackage{amssymb}
\usepackage{amsmath}
\usepackage[left=2cm,right=2cm,top=2cm,bottom=2cm]{geometry}
\usepackage{multirow}
\usepackage[dvipsnames]{xcolor}
\usepackage{gensymb} %symbole degré

\usepackage[T1]{fontenc}
\usepackage{fouriernc}
\usepackage{graphicx}

\usepackage{setspace}
\singlespacing

\usepackage{pstricks}
\usepackage{xkeyval}
\usepackage{pst-labo}
\usepackage{chemist} %chimie
\usepackage[version=4]{mhchem}%chimie

\usepackage{tcolorbox}
\usepackage{multicol}
\usepackage{caption}
\usepackage{fancyhdr}
\pagestyle{fancy}
\renewcommand\headrulewidth{1pt}
\fancyhead[L]{Lycée Denis Diderot }
\fancyhead[R]{Classe de 1 Bac Pro }
\setcounter{page}{1}\fancyfoot[C]{page \thepage}
\fancyfoot[R]{Année 2024-2025}
\fancyfoot[L]{Alexis RUIZ}
\usepackage{multirow,tabularx,booktabs}
\usepackage[tikz]{bclogo} 

\newcommand{\cadreblanc}[1]{%
\medskip\framebox[\linewidth][c]{%
\rule{0mm}{#1mm}}\par
\medskip}

\usepackage{awesomebox}
\setlength{\aweboxvskip}{0mm}
\usepackage{pifont}
\usepackage{chemmacros}
\chemsetup{modules=all}
\setlength{\columnseprule}{.25mm}


\newcommand*\acocher{\quitvmode{\fboxsep0pt \fboxrule0.8pt \fbox{\vrule height1.8ex depth.1ex width0pt \vrule height0pt depth0pt width1.9ex }}}

\newcolumntype{C}[1]{>{\centering\arraybackslash }b{#1}}
\renewcommand{\arraystretch}{1.3}
\newcommand{\code}{\awesomebox[black]{2pt}{

\includegraphics[scale=0.15]{Chapitre 1 - Energie et puissance électrique/images/frame.png}
}{black}}

\newcommand{\moteur}{\awesomebox[black]{2pt}{

\includegraphics[scale=1]{images/moteur.jpg} 
}{black}}

\newcommand{\defconclusion}{\awesomebox[black]{2pt}{\rotatebox{90}{\Large{\textbf{Conclusion}}}}{black}}

\newcommand{\defbox}{\awesomebox[black]{2pt}{
\includegraphics[scale=0.5]{Chapitre 1 - Energie et puissance électrique/images/appel exam.png}}{black}}

\newcommand{\docbox}{\awesomebox[black]{2pt}{
\includegraphics[scale=0.1]{Chapitre 4 - Solutions aqueuses/images/livre.png}}{black}}

  \newenvironment{itemize*}%  % listes compactes
    {\begin{termite}%
      \setlength{\itemsep}{0pt}%
      \setlength{\parskip}{0pt}}%
    {\end{itemize}}

\frenchbsetup{StandardLists=true}


%%%%%%%%%%%%%%%%%%%%%%%%%%%%%%%%%%%%%%%%%%%%%%%%%%%%%%%%%%
%                                                        %
%              Evaluation par compétences                %
%                                                        %
%%%%%%%%%%%%%%%%%%%%%%%%%%%%%%%%%%%%%%%%%%%%%%%%%%%%%%%%%%



\newcommand\SA{\marginpar{\scalebox{.8}{\framebox[1.6ex]{\rule{0ex}{1.6ex}\raisebox{.75ex}{\rule{1.6ex}{.5pt}}}\hspace{-.3pt}\framebox[3.2ex]{S}}}}
\newcommand\A{\marginpar{\scalebox{.8}{\framebox[1.6ex]{\rule{0ex}{1.6ex}\raisebox{.75ex}{\rule{1.6ex}{.5pt}}}\hspace{-.3pt}\framebox[3.2ex]{A}}}}
\newcommand\R{\marginpar{\scalebox{.8}{\framebox[1.6ex]{\rule{0ex}{1.6ex}\raisebox{.75ex}{\rule{1.6ex}{.5pt}}}\hspace{-.3pt}\framebox[3.2ex]{R}}}}
\newcommand\V{\marginpar{\scalebox{.8}{\framebox[1.6ex]{\rule{0ex}{1.6ex}\raisebox{.75ex}{\rule{1.6ex}{.5pt}}}\hspace{-.3pt}\framebox[3.2ex]{V}}}}
\newcommand\C{\marginpar{\scalebox{.8}{\framebox[1.6ex]{\rule{0ex}{1.6ex}\raisebox{.75ex}{\rule{1.6ex}{.5pt}}}\hspace{-.3pt}\framebox[3.2ex]{C}}}}


%%%%%%%%%%%%%%%%%%%%%%%%%%%%%%%%%%%

\begin{document}
%%%%%%%%%%%%%%%%%%%%%%%%%

\flushleft \textbf{NOM : \hfill Classe : \hspace{3cm} ~~\\[2mm]
Prénom:\\[0,4cm]
}
\hrule\
\\[0,2cm]

\begin{center}
 \LARGE{TP4 - L'eau pour le nourrisson}
\end{center}

\vspace{2  mm}

\normalsize

\hrule\
\\[0,2cm]
   \begin{center}
   \textbf{Le soin et la rédaction interviendront dans la notation.}
   \end{center} 
%%%%%%%%%%%%%%%%%%%%%%%%%%%
\noindent%
\begin{center}
    

\begin{minipage}{.95\textwidth}
Une eau peut être consommée en toute sécurité si ses caractéristiques microbiennes, chimiques et physiques répondent à certaines normes. \\[2mm]
Pour les nourrissons, la source d'eau est quasi-exclusivement l'eau contenue dans le lait qu'on lui prépare. \\

\end{minipage}%
\end{center}

\vfill 

\newtcolorbox{mybox}{colback=white!10!white,
colframe=black!5!black}
\begin{mybox}
\begin{center}    
\textbf{\large L'eau de Vittel utilisée pour préparer le biberon du nourrisson présente-t-elle des risques pour sa santé ? }\end{center}
\end{mybox}

\vfill 

  \docbox{\fbox{
\begin{minipage}{14cm}{Pour répondre aux questions, vous pourrez vous aider des documents présents dans le dossier documentaire.
  }
\end{minipage}
}}

\vfill 


\begin{enumerate}
    \item Nommer les critères chimiques de potabilité d'une eau. \SA  \\[5mm]

    \dotfill \\[3mm]

    \dotfill \\[3mm]

    \dotfill \\[3mm]

    \dotfill \\[3mm]

    

    

    \item Indiquer les risques sur la santé d'un nourrisson d'un apport trop important en ions sulfate. \SA \\[5mm]

    \dotfill \\[3mm]

    \dotfill \\[3mm]

    \item Relever la concentration massique maximale, en milligramme par litre,  en ions sulfate préconisée pour la réalisation d'un biberon. \SA  \\[5mm]

    \dotfill \\[3mm]

    \dotfill \\[3mm]





\vfill 

\newpage

\item Afin de connaître la concentration massique en ions sulfates  de l'eau de Vittel, on dose, par conductimétrie, un échantillon de cette eau. \\
Le dispositif expérimental est représenté ci-dessous. Compléter la description du protocole expérimental en précisant les différentes étapes . \\
Vous vous aiderez de la proposition de protocole en page 4.   \A
\vspace{5mm}





   
\noindent%
\begin{minipage}{.5\textwidth}%
\fbox{\includegraphics[scale=0.4]{Chapitre 4 - Solutions aqueuses/images/14 - schéma legendé.png}}
\end{minipage}%
\hfill
\begin{minipage}{.43\textwidth}%
\begin{center}
\fbox{Description du protocole}\\[5mm]

    \dotfill \\[3mm]

    \dotfill \\[3mm]
    
    \dotfill \\[3mm]

    \dotfill \\[3mm]
    
    \dotfill \\[3mm]

    \dotfill \\[3mm]
    
    \dotfill \\[3mm]

    \dotfill \\[3mm]

    
    \dotfill \\[3mm]

    \dotfill \\[3mm]
    
    \dotfill \\[3mm]

    \dotfill \\[3mm]

    
\end{center}\end{minipage}%



\vfill 

  \defbox{\fbox{
\begin{minipage}{13cm}{\textbf{Appel n\degree 1} : Appeler M. Ruiz afin de lui présenter oralement le protocole }
  
  \end{minipage}}}  \C

  \vfill 

  \item Réaliser les trois premières étapes du protocole à l'aide de la fiche technique "ExAO dosage".   \R

  \vfill 
  \defbox{\fbox{
\begin{minipage}{13cm}{\textbf{Appel n\degree 2} : Appeler M. Ruiz afin de réaliser devant lui les premières mesures de l'étape 4 }  \end{minipage}}}      \R

\vfill

\item Déterminer \chemform{V_{EQ}} à l'équivalence à l'aide de la fiche technique "ExAO volume équivalent". 
\begin{center}
    \fbox{\begin{minipage}{7cm}
    \vspace{3mm}
    \begin{center}
           \LARGE{\chemform{V_{EQ}} = ................ mL}
    \end{center}
    \end{minipage}
    }
\end{center}

\vfill 


\newpage


\item En utilisant la formule \large{$C_m\left(\ce{{SO_4^{2}-}}\right)=\dfrac{C_2 \times V_{EQ}}{V_1} \times M\left(\ce{SO_4^{2}-}\right)$}

\normalsize
, calculer la valeur de la concentration massique en ions sulfate de l'échantillon d'eau du robinet sachant que : 


\begin{itemize}
    

    \item $C_2$ : concentration molaire de la solution de chlorure de baryum
    \item $V_1$ : volume de l'échantillon d'eau. 
    \item $M(\ce{SO_4^{2}-})$ : Masse molaire des ions sulfate, vous utiliserez le tableau périodique pour la calculer. 
\end{itemize}
    \cadreblanc{30}  \R \SA

    \item Répondre à la problématique, justifier votre réponse :   

        
    \dotfill \\[3mm]

    \dotfill \\[3mm]
    
    \dotfill \\[3mm]

    \dotfill \\[3mm]  \V \C

\vfill 


\begin{center}
    

\begin{mybox}
\begin{center}    
\textbf{\Large {Dossier documentaire}}\end{center}
\end{mybox}
\end{center}

\normalsize
\textbf{Document 1 : Critères de potabilité de l'eau. } \\
Pour pouvoir être considérée comme potable en toute sécurité, une eau du robinet doit vérifier plusieurs critères de potabilité. Certains de ces critères sont indiqués ci-dessous : 
\begin{center}
    

\begin{tabular}{|C{8cm}|C{4cm}|}
\hline
Critères de potabilité & Valeur limite \\
\hline
Concentration massique en ions sulfate & < 250 mg/L \\
\hline
Concentration massique en ions chlorure & < 200 mg/L \\
\hline
Concentration massique en ions potassium & < 12 mg/L \\
\hline
pH  & 6,5 < pH < 9 \\
\hline
TH soit la dureté de l'eau & > 15 degrés français \\
\hline



\end{tabular}
\end{center}

\vfill 


\textbf{Document 2 : } \\

Les reins des nourrissons sont immatures jusqu'à l'âge d'un an, une surcharge en minéraux peut avoir des conséquences néfastes. Par exemple, un apport trop important d'ions sulfate, de formule \ce{SO_4^{2}-}, est nuisible pour les nourrissons car il peut provoquer des diarrhées et réduire l'absorption du calcium. 

Pour la réalisation des biberons des nourrissons, il est recommandé d'utiliser une eau dont la concentration massique en ions sulfate est inférieure à 140 mg/L. 

\vfill 

\newpage

\begin{center}
    

\begin{mybox}
\begin{center}    
\textbf{\Large {Proposition de protocole}}\end{center}
\end{mybox}

\fbox{
\begin{minipage}{15cm}
\begin{center}
    \includegraphics[scale=0.5]{Chapitre 4 - Solutions aqueuses/images/14 - schéma legendé.png}
\end{center}
    
\end{minipage}
}

\end{center}
\vfill 
\begin{center}
    

\fbox{

\begin{minipage}{15cm}
 \vspace{1mm}
    \textbf{Étape 1 : }
    \begin{itemize}
        \item Remplir la burette avec la solution de chlorure de Baryum.  
        \item  Ajuster le zéro. 
        \end{itemize}

    \textbf{Étape 2 : }
\begin{itemize}
    \item Introduire dans le bécher "DOSAGE" 100mL de la solution étiqueté A, on mesurera à l'aide de l'éprouvette graduée. 

    \item Introduire un barreau aimanté dans le bécher et le placer sur l'agitateur magnétique, sous la burette graduée. 
    
    \end{itemize}

   \textbf{ Étape 3  : }
    \begin{itemize}
        \item Positionner le capteur conductimétrique dans le bécher et paramétrer l'acquisition. 
    \end{itemize}

    \textbf{Etape 4 : }
    \begin{itemize}
        \item Lancer l'acquisition 
        \item Verser la solution de chlorure de Baryum 2mL par 2 mL. 
    \end{itemize}

    \vspace{2mm}

    
\end{minipage}
}
\end{center}

\vfill 
\end{enumerate}

\end{document}
